(1) 单模热平衡中数态权重正比于 \(q^n\)，其中 \(q=e^{-\hbar\omega/(k_BT)}\in(0,1)\)。正规化配分函数为 \(Z=\sum_{n\ge0}q^n=(1-q)^{-1}\)，所以 \(P_n=(1-q)q^n\)。由 \(\bar n=q/(1-q)\) 解出 \(q=\bar n/(1+\bar n)\)，故 \(P_n=\bar n^n/(1+\bar n)^{n+1}\)。
(2) 母函数 \(G(z)=\sum_nP_n z^n=(1-q)/(1-qz)\)，于是 \(G'(1)=\bar n\)、\(G''(1)=2\bar n^2\)。据 \(\operatorname{Var}(n)=G''(1)+G'(1)-G'(1)^2\)，得方差 \(\bar n(1+\bar n)\)；Fano 因子 \(F=\operatorname{Var}(n)/\bar n=1+\bar n\)，二阶关联 \(g^{(2)}(0)=\langle n(n-1)\rangle/\bar n^2=2\)，前提是 \(\bar n>0\)。
(3) 同平均的相干态具有泊松分布 \(P_n^{\rm coh}=e^{-\bar n}\bar n^n/n!\)，故 \(\operatorname{Var}_{\rm coh}(n)=\bar n\)、\(F_{\rm coh}=1\)、\(g_{\rm coh}^{(2)}(0)=1\)。两者的真空概率分别为 \(P_0^{\rm th}=1/(1+\bar n)\) 与 \(P_0^{\rm coh}=e^{-\bar n}\)；热态的零光子概率较高，同时高光子数尾部也较重。
(4) 热态 \(g^{(2)}(0)=2\) 表示两次接近同时的光子计数存在正相关，称为聚束；相干态的 \(g^{(2)}(0)=1\) 是泊松无聚束基准。严格真空 \(\bar n=0\) 时比值定义的分母为零，不能直接宣称真空的可测 \(g^{(2)}\) 等于某常数；应讨论非零占据的极限。光子能量量纲由 \(\hbar\omega\) 提供，\(\bar n\)、\(F\) 与 \(g^{(2)}\) 均无量纲。
